% ================================================================
% RESEARCH MEMO: Operadic Characterization of Chiral Koszulness
% Agent: RED-2 (operadic homotopy theorist)
% Date: 2026-03-17
% ================================================================

\documentclass[11pt]{article}
\usepackage{amsmath,amssymb,amsthm}
\usepackage[margin=1.2in]{geometry}
\usepackage{enumitem}

\newtheorem{conjecture}{Conjecture}
\newtheorem{question}{Question}
\newtheorem{observation}{Observation}
\theoremstyle{definition}
\newtheorem{proposal}{Proposal}
\theoremstyle{remark}
\newtheorem*{assessment}{Utility Assessment}
\newtheorem*{obstacle}{Obstacles}
\newtheorem*{literature}{Literature}

\newcommand{\Pch}{\mathcal{P}^{\mathrm{ch}}}
\newcommand{\Mbar}{\overline{\mathcal{M}}}
\newcommand{\FM}{\overline{C}}
\newcommand{\Def}{\mathrm{Def}}
\newcommand{\MC}{\mathrm{MC}}
\newcommand{\Sh}{\mathrm{Sh}}
\newcommand{\Heis}{\mathcal{H}}
\newcommand{\Vir}{\mathrm{Vir}}
\newcommand{\cA}{\mathcal{A}}

\title{%
  Operadic Characterization of Chiral Koszulness:\\
  Five Directions and the Shadow Tower Termination Problem}
\author{RED-2 (Operadic Homotopy Theory)}
\date{17 March 2026}

\begin{document}
\maketitle

\section*{Executive Summary}

The monograph's current characterization of chiral Koszulness
(Theorem~8.8 = \texttt{thm:pbw-koszulness-criterion}) reduces to
classical Koszulness of the PBW-associated graded.  This is
effective---it handles Heisenberg, Kac--Moody, Virasoro, and
Yangians---but it is a \emph{detection criterion}, not a
\emph{structural explanation}.  It says nothing about \emph{why}
certain algebras are Koszul, only \emph{that} they are.

The shadow Postnikov tower termination pattern
(Heisenberg at arity~2, affine at~3, $\beta\gamma$ at~4,
Virasoro at~$\infty$) is the most striking unexplained structure in
the monograph.  Below I propose five directions for an operadic
explanation.  My main thesis is:

\medskip
\noindent\textbf{The termination arity $r_{\max}(\cA)$ measures the
\emph{operadic complexity} of $\cA$ as an algebra over the chiral
operad $\Pch$, specifically the arity at which the minimal
$A_\infty$-model of $\cA$ becomes intrinsically formal.}

% ================================================================
\section{Direction A: Operadic vs.\ Algebra-Level Koszulness}
% ================================================================

\subsection{What the monograph currently says}

The monograph operates at two levels of Koszul duality
simultaneously (Remark~\texttt{rem:operadic-vs-algebraic-kd} in
\texttt{algebraic\_foundations.tex}, lines 1108--1133):

\begin{enumerate}[label=(\alph*)]
\item \emph{Operadic level}: $\mathrm{Com}^! = \mathrm{Lie}$,
  $\mathrm{Ass}^! = \mathrm{Ass} \otimes \mathrm{sgn}$.
  The chiral operad $\Pch$ is identified with
  $\Omega^*_{\log}(\FM_n(X))$ (Theorem~\texttt{thm:geometric-bridge},
  line~1169), and its Koszul dual cooperad is
  $(\Pch)^{!\prime}(n) \cong H_*(\FM_n(X))$
  (line~1184).
\item \emph{Algebraic level}: a chiral algebra $\cA$ is
  \emph{chiral Koszul} if the PBW-associated graded
  $\mathrm{gr}_F \cA$ is classically Koszul
  (Definition~\texttt{def:chiral-koszul-morphism}, the PBW criterion
  Theorem~\texttt{thm:pbw-koszulness-criterion}).
\end{enumerate}

The $E_1$-chiral Koszul duality theorem
(\texttt{thm:e1-chiral-koszul-duality}, line~2166) explicitly
derives the algebra-level bar-cobar equivalence from the
\emph{operadic} Koszulness of $\mathrm{Ass}$ via the transfer lemma
(\texttt{lem:operadic-koszul-transfer}, line~2268).  This is the
right framework.  The question is whether the analogous story works
for the full $E_\infty$-chiral (vertex algebra) case.

\subsection{Precise conjecture}

\begin{conjecture}[Operadic inheritance of chiral Koszulness]
\label{conj:operadic-inheritance}
Let $\Pch$ be the chiral operad on a smooth curve $X$, and let
$\cA$ be a $\Pch$-algebra (vertex algebra).
\begin{enumerate}
\item $\Pch$ is Koszul as an operad internal to
  $\mathcal{D}\textnormal{-mod}(X)$ equipped with the chiral tensor
  product.
\item Every $\Pch$-algebra $\cA$ that is PBW-Koszul in the sense of
  \texttt{thm:pbw-koszulness-criterion} is automatically Koszul
  \emph{as an algebra over a Koszul operad}, in the sense that the
  operadic bar-cobar adjunction
  $B_{\Pch} \dashv \Omega_{\Pch}$ restricts to an equivalence on
  the full subcategory of PBW-Koszul algebras.
\item Conversely, if $\Pch$ is Koszul and $\cA$ is
  $\Pch$-free on generators $V$ modulo quadratic-to-PBW relations,
  then $\cA$ is chiral Koszul.
\end{enumerate}
\end{conjecture}

\begin{obstacle}
The chiral operad $\Pch$ is not quadratic in the classical sense.
Its spaces $\Pch(n) = \Omega^*_{\log}(\FM_n(X))$ are
infinite-dimensional and carry mixed Hodge structures.  Koszulness
of $\Pch$ would need to be proved in the enriched sense
(internal to $\mathcal{D}$-mod$(X)$), following the transfer
strategy already used for $\mathrm{Ass}$ in
Lemma~\texttt{lem:operadic-koszul-transfer}.  The key subtlety:
$\Pch$ is \emph{not} the base-change of a classical operad.  It is
intrinsically geometric.  One needs the formality of the
FM compactification (Kontsevich) to reduce to the combinatorial
operad $H^*(\FM_n)$.
\end{obstacle}

\begin{literature}
Francis--Gaitsgory \cite{FG12}: $\mathrm{Lie}^{\mathrm{ch}}
\dashv \mathrm{Com}^{\mathrm{ch}}$ adjunction on $\mathrm{Ran}(X)$.
Loday--Vallette \cite{LV12}, Theorem~11.4.1: bar-cobar for
algebras over Koszul operads.  Fresse \cite{Fresse-operads},
\S16: model-categorical bar-cobar for general symmetric monoidal
model categories.  The monograph's
Lemma~\texttt{lem:operadic-koszul-transfer} is the precise bridge.
\end{literature}

\begin{assessment}
\textbf{High utility, moderate difficulty.}  This conjecture would
unify the PBW criterion with the operadic framework and explain
\emph{why} the PBW strategy works: it works because $\Pch$ is
Koszul, and PBW-flat algebras over Koszul operads are automatically
Koszul.  The proof strategy is essentially already in the monograph
for $\mathrm{Ass}$; extending it to $\Pch$ requires the formality
input.
\end{assessment}

% ================================================================
\section{Direction B: Distributive Law Characterization}
% ================================================================

\subsection{The idea}

A vertex algebra OPE has the structure
\[
a(z) b(w) = \underbrace{\sum_{n \geq 0}
\frac{(a_{(n)} b)(w)}{(z-w)^{n+1}}}_{\text{singular part}}
+ \underbrace{:a(z)b(w):}_{\text{regular part}}.
\]
The singular part encodes the ``Lie-type'' data (poles, commutators),
and the regular part encodes the ``commutative-type'' data (normal
ordering).  The PBW theorem for vertex algebras says that the
multiplication map
\[
\mathrm{Sym}^{\mathrm{ch}}(V) \xrightarrow{\sim} \mathrm{gr}_F \cA
\]
is an isomorphism.  In the Loday--Vallette framework, PBW theorems
for algebras are consequences of \emph{distributive laws} between
operads: if $P = P_1 \circ_\lambda P_2$ (a distributive law
$\lambda: P_2 \circ P_1 \to P_1 \circ P_2$), then free
$P$-algebras have PBW bases indexed by $P_1$-monomials applied to
$P_2$-monomials.

\subsection{Precise conjecture}

\begin{conjecture}[Distributive law for chiral Koszulness]
\label{conj:distributive-law}
There exist sub-operads
$\Pch_{\mathrm{sing}} \hookrightarrow \Pch$ and
$\Pch_{\mathrm{reg}} \hookrightarrow \Pch$
encoding, respectively, the singular (OPE pole) and regular (normal
ordering) parts of the chiral product, together with a distributive
law
\[
\lambda: \Pch_{\mathrm{reg}} \circ \Pch_{\mathrm{sing}}
\longrightarrow
\Pch_{\mathrm{sing}} \circ \Pch_{\mathrm{reg}}
\]
such that:
\begin{enumerate}
\item $\Pch \cong \Pch_{\mathrm{sing}} \circ_\lambda
  \Pch_{\mathrm{reg}}$ as operads.
\item The PBW theorem for chiral algebras
  (the isomorphism $\mathrm{Sym}^{\mathrm{ch}}(V) \cong
  \mathrm{gr}_F \cA$) is a consequence of this distributive law.
\item A chiral algebra $\cA$ is Koszul if and only if both
  $\Pch_{\mathrm{sing}}$ and $\Pch_{\mathrm{reg}}$ are Koszul and
  the distributive law is compatible with the Koszul resolutions.
\end{enumerate}
\end{conjecture}

\begin{obstacle}
This is the most speculative direction.  The main difficulty is
that the singular/regular decomposition of the OPE is
\emph{not operadic} in a naive sense: the singular part at $n$
points involves pole orders that depend on the specific algebra
$\cA$ (e.g., double pole for Heisenberg, quartic pole for
Virasoro).  A genuine distributive law would need to operate at the
level of the chiral operad itself, not at the algebra level.

Moreover, Loday--Vallette's distributive law theory
(\cite{LV12}, \S8.6) applies to quadratic operads and produces
quadratic-linear algebras via the rewriting method.  The chiral
setting involves poles of \emph{all} orders, which goes beyond the
quadratic-linear framework.

A more realistic version: the distributive law exists at the level
of the \emph{PBW-graded} chiral operad
$\mathrm{gr}_F \Pch$, where it reduces to the classical
$\mathrm{Com} \circ \mathrm{Lie}$ distributive law (encoding the
PBW theorem for universal enveloping algebras).  The failure of
the distributive law to lift to the ungraded level is measured by
a sequence of obstructions, and the \emph{shadow tower} is
precisely this obstruction sequence.
\end{obstacle}

\begin{literature}
Loday--Vallette \cite{LV12}, \S8.6: distributive laws and PBW.
Dotsenko--Khoroshkin \cite{DK10}: Gr\"obner bases for operads,
rewriting approach to PBW.  Bremner--Dotsenko: algebraic operads
and rewriting systems.  For the vertex algebra PBW theorem: Li
\cite{Li96}, Frenkel--Ben-Zvi \cite{FBZ04}.
\end{literature}

\begin{assessment}
\textbf{Medium utility, high difficulty.}  If achieved, this would
give the \emph{deepest} structural explanation of chiral
Koszulness.  But the gap between the operadic distributive law
framework and the analytic structure of OPE poles is large.
The modified version (distributive law on the associated graded,
with the shadow tower as the obstruction to lifting) is more
realistic and connects directly to the monograph's existing
architecture.
\end{assessment}

% ================================================================
\section{Direction C: Modular Koszulness and Shadow Tower
Termination}
\label{sec:modular}
% ================================================================

\subsection{The key data}

From the monograph
(Example~\texttt{ex:tower-four-archetypes},
\texttt{higher\_genus\_modular\_koszul.tex} lines 8571--8597):

\begin{center}
\begin{tabular}{lcccc}
\hline
\textbf{Algebra} & $\kappa$ & $\mathfrak{C}$ & $\mathfrak{Q}$ &
$r_{\max}$ \\
\hline
Heisenberg $\Heis_k$ & $k$ & $0$ & $0$ & $2$ \\
Affine $\widehat{\mathfrak{sl}}_2$ & $\kappa$ &
$\kappa(x,[y,z])$ & $0$ (Jacobi) & $3$ \\
$\beta\gamma$ & $0$ & $0$ & $\mathrm{cyc}(m_3)$ & $4$ \\
Virasoro $\Vir_c$ & $c/2$ & $2x^3$ &
$\frac{10}{c(5c+22)}x^4$ & $\infty$ \\
\hline
\end{tabular}
\end{center}

The shadow Postnikov tower
(Definition~\texttt{def:shadow-postnikov-tower}) is a sequence of
MC truncations $\Theta_\cA^{\leq r}$ in the modular convolution
dg Lie algebra $\mathfrak{g}_\cA^{\mathrm{mod}}$, where
$\Theta_\cA^{\leq r}$ solves MC modulo weight $r+1$.  The
obstruction class $o_{r+1}(\cA) \in H^2(\cA^{\mathrm{sh}}_{r+1,0})$
controls extension.  Termination at $r_{\max}$ means
$\cA^{\mathrm{sh}}_{r,0} = 0$ for $r > r_{\max}$.

\subsection{Operadic explanation of termination}

\begin{conjecture}[Operadic complexity controls termination]
\label{conj:operadic-termination}
The shadow tower termination arity $r_{\max}(\cA)$ equals the
\emph{operadic complexity} of $\cA$, defined as the minimal arity
$r$ such that the transferred $A_\infty$ operations
$m_n^{\mathrm{tr}} = 0$ for all $n > r$ on the minimal model of
$\cA$ as a $\Pch$-algebra.
\end{conjecture}

Here is the proposed dictionary:

\begin{center}
\begin{tabular}{lll}
\hline
\textbf{Algebra} & \textbf{Minimal model} & $r_{\max}$ \\
\hline
Heisenberg & Strictly commutative ($m_n = 0$, $n \geq 3$) & $2$ \\
Affine KM & Strict Lie ($m_3 \neq 0$ from bracket, $m_n = 0$,
$n \geq 4$) & $3$ \\
$\beta\gamma$ & $A_\infty$ with $m_3$ and $m_4$ from contact
terms & $4$ \\
Virasoro & Full $A_\infty$ tower, $m_n \neq 0$ for all $n$ & $\infty$ \\
\hline
\end{tabular}
\end{center}

\textbf{Why this should be true.}  The shadow at arity $r$ is built
from graph sums (Definition~\texttt{def:modular-bar-hamiltonian},
line 8610) whose vertices carry the transferred operations
$\ell_{|v|}^{\mathrm{tr}}$ from the cyclic minimal model.  If the
minimal model has $m_n^{\mathrm{tr}} = 0$ for $n > r$, then all
graph sums with a vertex of valence $> r$ vanish, and the only
contributions to the shadow at arity $> r$ come from graphs whose
vertices all have valence $\leq r$.  The obstruction class
$o_{r+1}$ involves the bracket
$[\Theta^{\leq r}, \Theta^{\leq r}]_{r+1}$, which is a sum over
graphs with two vertices sewn together.  If the vertex operations
truncate at arity $r$, and the bracket produces arity $r+1$ only
by combining vertices of arity $\leq r$, then $o_{r+1} = 0$ if and
only if the ``Jacobi-type'' identity for these lower-arity
operations holds.

The pattern:
\begin{itemize}
\item \textbf{Heisenberg}: $m_n^{\mathrm{tr}} = 0$ for $n \geq 3$
(abelian).  All graph sums with $\geq 3$-valent vertices vanish.
The modular content is purely from edge contractions (propagators),
giving the Gaussian $\kappa \cdot \lambda_g$ formula.

\item \textbf{Affine}: $m_2^{\mathrm{tr}}$ is the Lie bracket,
$m_3^{\mathrm{tr}} = 0$ by transfer from the Koszul resolution (the
$A_\infty$ structure on $\mathrm{Ext}^*(\mathbb{C}, \mathbb{C})$
for $U(\mathfrak{g})$ has $m_3 = 0$ because the Koszul resolution
is \emph{minimal} and the Chevalley--Eilenberg complex is formal
through degree 3).  The cubic shadow $\mathfrak{C}$ comes from
$m_2$ composed with itself (the $s$-$t$-$u$ channels), and the
quartic obstruction vanishes by the Jacobi identity.

\item \textbf{$\beta\gamma$}: The normal-ordered product
$:\beta\gamma:$ gives a nontrivial $m_3^{\mathrm{tr}}$ (the contact
term from the quartic singularity in the $\beta\gamma$ OPE after
subtracting the simple pole), but $m_4^{\mathrm{tr}}$ is exact
(killed by rank-one abelian rigidity,
\texttt{thm:nms-rank-one-rigidity}).

\item \textbf{Virasoro}: The composite field $:TT:$ and its
descendants generate an infinite sequence of nonvanishing
transferred operations.  The $A_\infty$ minimal model of $\Vir_c$
is genuinely infinite: $m_n^{\mathrm{tr}} \neq 0$ for all $n$
because the OPE involves a quartic pole, composite fields of
arbitrarily high weight, and the central charge enters nonlinearly
at every order (as noted in
\texttt{chiral\_koszul\_pairs.tex}, line~2896).  The quintic is
forced because $\{\mathfrak{C}, \mathfrak{Q}\}_H \neq 0$
(Theorem~\texttt{thm:nms-virasoro-quintic-forced}).
\end{itemize}

\begin{obstacle}
The precise relationship between $m_n^{\mathrm{tr}}$ and the shadow
obstruction $o_n$ is not a direct equality.  The obstruction
$o_{r+1}$ involves \emph{all} graphs of weight $r+1$, including
graphs with low-valence vertices connected by many edges (loops).
A vertex with $m_k^{\mathrm{tr}} = 0$ for $k > r$ does not
immediately imply $o_{r+1} = 0$, because multi-loop graphs with
vertices of valence $\leq r$ can contribute to weight $r+1$.

The precise statement would need: if $m_k^{\mathrm{tr}} = 0$ for
$k > r$ AND the genus-loop operator $\Lambda_P$ acting on
arity-$(r+2)$ shadows from lower-arity vertices gives zero, THEN
$o_{r+1} = 0$.  The genus-loop condition is an additional
constraint beyond the truncation of the minimal model.
\end{obstacle}

\begin{literature}
Merkulov \cite{Mer10}: minimal models and $A_\infty$ transfer for
operadic algebras.  Kadeishvili \cite{Kad80}: $A_\infty$ minimal
models.  Kontsevich--Soibelman \cite{KS06}: formal moduli from
$L_\infty$ and $A_\infty$.  The graph sum formula
(\texttt{def:modular-bar-hamiltonian}) is the Feynman expansion of
the quantum $L_\infty$ structure
(Chuang--Lazarev \cite{CL10}).
\end{literature}

\begin{assessment}
\textbf{Very high utility, moderate difficulty.}  This is the most
promising direction.  It gives a precise operadic meaning to the
shadow termination pattern and connects it to the classical theory
of $A_\infty$ minimal models.  The main work is making the
graph-sum argument rigorous: showing that truncation of the
minimal model controls termination of the shadow tower, accounting
for the genus-loop contributions.

This could be stated as a theorem in the monograph, not merely a
conjecture, if the graph-sum analysis is carried out carefully.
\end{assessment}

% ================================================================
\section{Direction D: Koszul Duality of the Chiral Operad}
% ================================================================

\subsection{What is $(\Pch)^!$?}

The monograph identifies $\Pch(n) \cong \Omega^*_{\log}(\FM_n(X))$
(\texttt{algebraic\_foundations.tex}, line~1171) and the Koszul dual
cooperad $(\Pch)^{!\prime}(n) \cong H_*(\FM_n(X))$ (line~1184).
Since $H^*(\FM_n(\mathbb{C}))$ is generated by the Arnold classes
$\omega_{ij}$ with the Arnold relations, and the homology
$H_*(\FM_n(\mathbb{C}))$ is generated by the dual classes
$\alpha_{ij}$ with the dual relations, the Koszul dual operad
$(\Pch)^!$ is:

\[
(\Pch)^!(n) = H_*(\FM_n(X), \mathbb{C})
= \text{the Lie-type operad generated by ``collision classes''.}
\]

On $\mathbb{P}^1$:
$H^*(\FM_n(\mathbb{P}^1)) \cong \mathrm{Com}(n)$ (the commutative
operad), and
$H_*(\FM_n(\mathbb{P}^1)) \cong \mathrm{Lie}(n)$ (by the partition
lattice theorem, \texttt{algebraic\_foundations.tex},
Theorem~\texttt{thm:com-lie}, line~1236).  So:
\[
(\Pch)^! \Big|_{X = \mathbb{P}^1}
\cong \mathrm{Lie}^{\mathrm{ch}}.
\]
This recovers the Francis--Gaitsgory
$\mathrm{Lie}^{\mathrm{ch}} \dashv \mathrm{Com}^{\mathrm{ch}}$
adjunction.

\subsection{Precise question}

\begin{question}
For a higher-genus curve $X$ (or the universal curve over
$\Mbar_{g,n}$), what is $(\Pch_X)^!$?  Does it depend on the
complex structure of $X$?  Is the Koszul dual cooperad
$(\Pch_X)^{!\prime}(n) = H_*(\FM_n(X))$ formal
(i.e., does the cooperad structure on chains $C_*(\FM_n(X))$
transfer to cohomology)?
\end{question}

At genus $g \geq 1$, $H^1(\FM_n(X)) \neq 0$ (from the homology of
the curve itself), and the chiral operad $\Pch_X$ is \emph{not}
the base-change of $\mathrm{Com}$.  The Koszul dual
$(\Pch_X)^!$ would be a \emph{genus-dependent} Lie-type operad,
whose arity-$n$ spaces encode the interaction between the OPE and
the topology of $X$.

\begin{conjecture}[Genus-dependent Koszul dual]
\label{conj:genus-koszul-dual}
For a smooth curve $X$ of genus $g$:
\begin{enumerate}
\item The cooperad $\{C_*(\FM_n(X))\}$ is formal if and only if
$g = 0$.
\item At genus $g \geq 1$, the non-formality is controlled by the
period matrix $\Omega_g \in \mathfrak{H}_g$, and the resulting
``Lie-type'' operad $(\Pch_X)^!$ carries a connection over the
Siegel upper half-space (equivalently, over $\Mbar_g$).
\item The curvature of this connection is the modular characteristic
$\kappa(\cA)$: the obstruction to $(\Pch_X)^!$ being independent
of the complex structure.
\end{enumerate}
\end{conjecture}

This is essentially saying: the genus-$g$ bar complex
$\bar{B}^{(g)}(\cA)$ computes the $(\Pch_X)^!$-coalgebra
structure on $\cA$, and the curvature
$d_{\mathrm{fib}}^2 = \kappa \cdot \omega_g$ is the failure of
this coalgebra structure to be flat over $\Mbar_g$.

\begin{obstacle}
The formality of configuration spaces on higher-genus curves is a
deep open question.  On $\mathbb{C}$ (or $\mathbb{P}^1$),
Kontsevich formality gives $C^*(\FM_n(\mathbb{C})) \simeq
H^*(\FM_n(\mathbb{C}))$ as operads.  On $\Sigma_g$ for $g \geq 1$,
this fails: $C^*(\mathrm{Conf}_n(\Sigma_g))$ is \emph{not} formal
in general (there are nontrivial Massey products from the
interaction of the curve's homology with the configuration space
structure).  This non-formality is closely related to the
non-vanishing of the shadow tower at higher arities.
\end{obstacle}

\begin{assessment}
\textbf{High utility, high difficulty.}  This direction connects
the shadow tower directly to the homotopy theory of configuration
spaces on curves, which is an active area of research.  The
payoff would be enormous: a geometric explanation of $\kappa$ as
the curvature of a connection on the Koszul dual operad.  But the
technical machinery (formality on curves, connections on operads
over moduli) is demanding.
\end{assessment}

% ================================================================
\section{Direction E: Transfer, Formality, and Operadic Koszulness}
% ================================================================

\subsection{The framework}

The monograph states that bar-cobar preserves quasi-isomorphisms
``because it is a quantum $L_\infty$ functor''
(\texttt{CLAUDE.md}).  More precisely: the convolution
$L_\infty$-algebra $\mathrm{hom}_\alpha(C, A)$ of
Robert-Nicoud--Wierstra carries a transferred $L_\infty$ structure,
and the bar-cobar adjunction is the MC moduli of this structure
(Theorem~\texttt{thm:modular-quantum-linfty}, line~7537).

The homotopy transfer theorem (HTT) says: given a deformation
retract $V \hookrightarrow A \twoheadrightarrow V$ with
$V = H^*(A)$, the $A_\infty$ structure on $A$ transfers to a
\emph{minimal} $A_\infty$ structure on $V$.  The transferred
operations $m_n^{\mathrm{tr}}$ are sums over planar trees with
$n$ leaves, each internal edge decorated by the contracting
homotopy $h$ and each vertex by the original operations.

\begin{conjecture}[Formality characterization of Koszulness]
\label{conj:formality-koszulness}
A chiral algebra $\cA$ is Koszul if and only if the convolution
$L_\infty$-algebra
$\mathrm{hom}_\alpha(\{C_*(\FM_n(X))\},
\mathrm{End}_\cA)$
is \emph{formal}: i.e., quasi-isomorphic to its cohomology
(a dg Lie algebra with zero differential) as an $L_\infty$-algebra.

Moreover:
\begin{enumerate}
\item $\cA$ is Koszul with $r_{\max} = 2$
  (shadow tower terminates at arity~2) if and only if the
  convolution $L_\infty$-algebra is \emph{intrinsically formal}
  (all higher brackets $\ell_k = 0$ for $k \geq 3$ on cohomology).
\item $r_{\max} = r$ if and only if the convolution
  $L_\infty$-algebra is \emph{formal through arity $r$} but not
  through arity $r+1$: the transferred brackets $\ell_k^{\mathrm{tr}}$
  vanish for $k > r$ but $\ell_{r+1}^{\mathrm{tr}}$ represents a
  nontrivial class in the Harrison cohomology of the $L_\infty$
  deformation complex.
\item $r_{\max} = \infty$ (non-termination, e.g., Virasoro) if and
  only if the convolution $L_\infty$-algebra is \emph{not formal at
  any finite level}: there exist nontrivial transferred brackets at
  every arity.
\end{enumerate}
\end{conjecture}

\begin{observation}[The formality obstruction tower IS the shadow
tower]
The obstructions to formality of an $L_\infty$-algebra are organized
by a Postnikov-type tower: at each stage, one tries to kill the
lowest-arity nontrivial transferred bracket by a quasi-isomorphism.
The obstruction to killing $\ell_n^{\mathrm{tr}}$ lives in a
Harrison-type cohomology group.  This tower is formally identical
to the shadow Postnikov tower:
\begin{itemize}
\item $\Theta^{\leq 2} = \kappa$ $\leftrightarrow$ binary bracket
  (dg Lie structure).
\item $o_3 = $ cubic shadow obstruction $\leftrightarrow$
  $\ell_3^{\mathrm{tr}}$, the ternary Massey-type product.
\item $o_4 = $ quartic resonance $\leftrightarrow$
  $\ell_4^{\mathrm{tr}}$, the quaternary obstruction.
\end{itemize}
\end{observation}

This is consistent with the monograph's explicit formula for
$\ell_3^{(0)}$
(Construction~\texttt{constr:explicit-convolution-linfty},
line~7582): it is the sum over the three binary trees with
3~leaves ($s$, $t$, $u$ channels), each internal edge contracted
with the homotopy propagator.  On cohomology,
$[\ell_3^{(0)}] = 0$ in $H^*(\Mbar_{0,4}) = H^*(\mathbb{P}^1)$
because $H^1(\mathbb{P}^1) = 0$.  But at the chain level,
$\ell_3^{(0)}$ is nontrivial and feeds the cubic shadow.

The quaternary bracket $\ell_4^{(0)}$ involves 15 trivalent trees
with 4 leaves on $\Mbar_{0,5}$.  Its cohomology class detects the
quartic shadow obstruction $o_4(\cA)$ (line~7644).

\begin{obstacle}
The identification of the formality obstruction tower with the
shadow Postnikov tower requires showing that the Harrison
cohomology controlling $L_\infty$ formality maps isomorphically to
the shadow algebra $\cA^{\mathrm{sh}}$.  This is plausible but not
automatic: the shadow algebra is defined as
$H_\bullet(\Def_{\mathrm{cyc}}^{\mathrm{mod}}(\cA))$
(Definition~\texttt{def:shadow-algebra}), which is the cyclic
deformation complex, not the Harrison complex.  The two coincide
for commutative algebras but differ in general.

Additionally, the genus contributions complicate the story:
$\ell_1^{(g)}$ (the BV operator at genus $g$) mixes with the
genus-0 formality obstructions.  A clean statement would need to
separate the genus-0 formality from the genus-$g$ corrections.
\end{obstacle}

\begin{assessment}
\textbf{Very high utility, moderate difficulty.}  This is the most
satisfying answer to the shadow tower termination question.  It
gives a dictionary:
\begin{center}
\begin{tabular}{ll}
\hline
\textbf{Shadow tower} & \textbf{Formality theory} \\
\hline
$\kappa$ (arity 2) & Binary Lie bracket \\
$\mathfrak{C}$ (arity 3) & Ternary Massey product \\
$\mathfrak{Q}$ (arity 4) & Quaternary formality obstruction \\
$r_{\max} < \infty$ & Finite-stage formality \\
$r_{\max} = \infty$ & Intrinsic non-formality \\
\hline
\end{tabular}
\end{center}
The Virasoro tower being infinite means: the convolution
$L_\infty$-algebra of $\Vir_c$ is \emph{intrinsically non-formal}.
The Heisenberg tower terminating at arity~2 means: the convolution
$L_\infty$-algebra is intrinsically formal (a genuine dg Lie
algebra, no higher brackets needed).  This is exactly the
content of the shadow archetype classification.
\end{assessment}

% ================================================================
\section{The Key Question: Operadic Explanation of the Termination
Pattern}
% ================================================================

Let me combine the threads.  The shadow tower termination pattern
\[
\Heis @2, \quad
\widehat{\mathfrak{g}} @3, \quad
\beta\gamma @4, \quad
\Vir @\infty
\]
admits the following unified operadic explanation, synthesizing
Directions C and E:

\begin{proposal}[The Operadic Complexity Theorem]
\label{prop:operadic-complexity}
For a chiral Koszul algebra $\cA$ over $\Pch$, define:
\begin{enumerate}
\item The \emph{$A_\infty$-depth} $d_\infty(\cA)$ = the largest
$n$ such that $m_n^{\mathrm{tr}} \neq 0$ in the minimal
$A_\infty$ model of $\cA$.
\item The \emph{$L_\infty$-formality level} $f_\infty(\cA)$ = the
largest $n$ such that $\ell_n^{(0),\mathrm{tr}} \neq 0$ in the
transferred quantum $L_\infty$ brackets on the convolution algebra.
\item The \emph{shadow depth} $r_{\max}(\cA)$ = the termination
arity of the shadow Postnikov tower.
\end{enumerate}

Then:
\[
r_{\max}(\cA) = d_\infty(\cA) = f_\infty(\cA).
\]
\end{proposal}

\textbf{Verification on examples:}
\begin{itemize}
\item \textbf{Heisenberg}: $d_\infty = 2$ (commutative, all
$m_n = 0$ for $n \geq 3$); $f_\infty = 2$ (convolution is a dg Lie
algebra, no higher brackets); $r_{\max} = 2$.
\checkmark

\item \textbf{Affine}: $d_\infty = 3$ (the Lie bracket gives
$m_2 \neq 0$; the Koszul resolution is minimal so $m_3 = 0$ on
the minimal model, but the cubic shadow
$\mathfrak{C} = \kappa(x,[y,z])$ comes from the \emph{convolution}
$\ell_3$ which is chain-level nontrivial even though
cohomologically trivial at genus 0---wait.  Actually,
$r_{\max} = 3$ means $o_4 = 0$, which is the Jacobi identity.
The $A_\infty$-depth is $2$ (binary Lie bracket), not $3$.
The cubic shadow comes from composing $m_2$ with itself, not from
$m_3$.  So we need $d_\infty(\cA)$ to count the ``effective arity''
including compositions.)

\textbf{Correction}: the right notion is not $m_n^{\mathrm{tr}}$
alone, but the \emph{effective operadic arity} = the arity at
which the graph sums from all vertices of the minimal model
can generate nontrivial contributions.  For affine: $m_2$ is
nonzero (the Lie bracket), and graphs with two binary vertices
sewn together give arity-3 contributions.  But graphs with
three binary vertices give arity-4 contributions, and $o_4 = 0$
because the Jacobi identity is an operadic relation.

So the refined statement is: $r_{\max}(\cA)$ is the arity at
which the \emph{operadic relations} of $\Pch$ kill all further
graph-sum contributions.  For $\mathrm{Lie}$: the Jacobi identity
is the unique relation in arity 3 of the Lie operad, and it kills
all arity-4 contributions.  This is $r_{\max} = 3$.

\item \textbf{$\beta\gamma$}: The $A_\infty$ minimal model has
$m_2 = 0$ (no Lie bracket: this is a free field), $m_3 \neq 0$
(from the normal-ordered product and the OPE contact terms).
Graphs with one ternary vertex and edges give contributions at
arities 3 and 4.  The quartic shadow $\mathfrak{Q}$ is nonzero,
but $o_5 = 0$ by rank-one abelian rigidity.
$r_{\max} = 4$.

\item \textbf{Virasoro}: $m_2 \neq 0$ (the Virasoro bracket),
the composite field $:TT:$ generates effective $m_4$-type
contributions, and the central charge prevents the graph sums from
terminating.  The obstruction
$\{C, Q\}_H \neq 0$ forces $\Sh_5 \neq 0$, and by
induction the tower is infinite.  $r_{\max} = \infty$.
\end{itemize}

\begin{conjecture}[Refined operadic complexity]
\label{conj:refined}
$r_{\max}(\cA)$ equals the \emph{relational arity} of $\cA$ as a
$\Pch$-algebra: the minimal $r$ such that all defining relations of
$\cA$ (including the operadic relations of $\Pch$ and the specific
OPE relations of $\cA$) can be expressed in terms of operations of
arity $\leq r$, AND the Jacobi-type higher syzygies generated by
these relations kill all obstruction classes $o_{r+1}, o_{r+2},
\ldots$ simultaneously.

This is the chiral analogue of the following classical fact: for
a quadratic algebra $A = T(V)/(R)$ with $R \subset V^{\otimes 2}$,
the Koszul property is equivalent to the quadratic relations $R$
generating all syzygies.  For a cubic algebra (relations in
$V^{\otimes 3}$), one needs $r = 3$.  The ``relational arity'' is
the operadic version of the relation degree.
\end{conjecture}

% ================================================================
\section{Synthesis and Recommendations}
% ================================================================

\subsection{Priority ranking}

\begin{enumerate}
\item \textbf{Direction E} (Formality characterization): highest
  priority.  The identification of the shadow tower with the
  formality obstruction tower for the convolution $L_\infty$-algebra
  is the cleanest answer to the termination question and is
  closest to being provable with existing technology.

\item \textbf{Direction C} (Operadic complexity): closely related
  to Direction~E.  The graph-sum analysis connecting
  $A_\infty$-depth to shadow termination should be developed in
  parallel.  This is the computational workhorse.

\item \textbf{Direction A} (Operadic inheritance): important for
  the foundational architecture.  Proving that $\Pch$ is Koszul
  (using Kontsevich formality) would ground the entire PBW framework
  in operadic theory.

\item \textbf{Direction D} (Genus-dependent Koszul dual): most
  ambitious, connecting to deep questions in algebraic geometry
  (formality of configuration spaces on curves).  Long-term.

\item \textbf{Direction B} (Distributive law): most speculative.
  The modified version (distributive law on associated graded, shadow
  tower as lifting obstruction) is viable but not yet precise enough
  to state as a theorem.
\end{enumerate}

\subsection{What the monograph should add}

I recommend the following additions:

\begin{enumerate}
\item A remark in \texttt{chiral\_koszul\_pairs.tex}, after
  Theorem~\texttt{thm:pbw-koszulness-criterion}, stating that the
  PBW criterion is a \emph{consequence} of the Koszulness of the
  chiral operad $\Pch$, in the same way that the $E_1$-chiral
  Koszul duality (Theorem~\texttt{thm:e1-chiral-koszul-duality}) is
  a consequence of the Koszulness of $\mathrm{Ass}$.

\item A conjecture in \texttt{higher\_genus\_modular\_koszul.tex},
  after Definition~\texttt{def:shadow-postnikov-tower}, stating the
  operadic complexity conjecture (Conjecture~\ref{conj:operadic-termination}):
  the shadow termination arity equals the $A_\infty$-depth of the
  minimal model.

\item A remark connecting the shadow Postnikov tower to the
  formality obstruction tower for the convolution $L_\infty$-algebra
  (Conjecture~\ref{conj:formality-koszulness}).  The key formula is
  already in the monograph
  (Construction~\texttt{constr:explicit-convolution-linfty}):
  the quantum $L_\infty$ brackets $\ell_n^{(g)}$ are the Feynman
  expansion, and their vanishing on cohomology is formality.
\end{enumerate}

\subsection{The bottom line}

The shadow tower termination pattern
$\Heis @2$, $\widehat{\mathfrak{g}} @3$, $\beta\gamma @4$,
$\Vir @\infty$ is not an accident.  It measures the
\emph{operadic complexity} of the chiral algebra: the depth of its
$A_\infty$ minimal model, equivalently, the formality level of
its modular convolution $L_\infty$-algebra.  Commutative algebras
(Heisenberg) are intrinsically formal.  Lie algebras (affine) have
one layer of non-formality, killed by the Jacobi identity.  Contact
algebras ($\beta\gamma$) have two layers, killed by abelian
rigidity.  The Virasoro algebra, with its quartic pole and infinite
family of composite fields, is intrinsically non-formal at every
level: its modular content is genuinely infinite.

This is the operadic explanation the monograph lacks.  The
technology to prove it is largely present: the graph-sum formula,
the explicit $L_\infty$ brackets, and the obstruction recursion.
What remains is the formal argument connecting these ingredients.

\end{document}
